\documentclass[12pt]{article}
\usepackage{amsmath, amssymb, amsthm}
\usepackage{geometry}
\usepackage{xcolor}
\geometry{margin=1in}

\title{Engineering Probability \\ Homework 1}
\author{Timothy Pollard, pollat@rpi.edu}
\date{September 6th 2025}

\begin{document}

\maketitle
\hrulefill

\section*{(1)}
An integrated circuit factory has three machines $X, Y$ and $Z$. Test one
integrated circuit produced by each machine. Either a circuit is acceptable ($a$) or it
fails ($f$). An observation is a sequence of three test results corresponding to the circuits
from machines $X, Y$ and $Z$ respectively. For example $aaf$ is the observation that the
circuits from $X$ and $Y$ pass the test and the circuit from $Z$ fails the test.

\begin{enumerate}
    \item[a] What are the elements of the sample space of this experiment?
    \textcolor{blue}{\item[] $S=\{aaa,aaf,afa,aff,faa,faf,ffa,fff\}$}
    \item[b] What are the elements of the sets:
    \begin{enumerate}
        \item[] $Z_F = \{\text{circuit from } Z \text{ fails}\} =$ \textcolor{blue}{$\{aaf, aff, faf, fff\}$}
        \item[] $X_A = \{\text{circuit from } X \text{ is acceptable}\} =$ \textcolor{blue}{$\{aaa, aaf, afa, aff\}$}
    \end{enumerate}
    \item[c] Are \(Z_F\) and \(X_A\) mutually exclusive? \textcolor{blue}{ No the element $aaf$ is in both sets.}
    \item[d] Are \(Z_F\) and \(X_A\) collectively exhaustive? 
    \item[]\textcolor{blue}{ No the element $ffa$ is not included in either of the sets.} 
    \item[e] What are the elements of the sets:
    \begin{enumerate}
        \item[] $C = \{\text{more than one circuit acceptable}\} = $ \textcolor{blue}{$\{aaa, aaf, afa, faa\}$}
        \item[] $D = \{\text{at least two circuits fail}\} =$ \textcolor{blue}{$\{aff, faf, ffa, fff\}$}
    \end{enumerate}
    \item[f] Are \(C\) and \(D\) mutually exclusive?
    \item[] \textcolor{blue}{Yes}
    \item[g] Are \(C\) and \(D\) collectively exhaustive?
    \item[]\textcolor{blue}{Yes}
\end{enumerate}

\section*{(2)}
There are two types of cellular phones, hand-held phones ($H$) that you carry
and mobile phones ($M$) that are mounted in vehicles. Phone calls can be classified by
the traveling speed of a user as fast ($F$) or slow ($W$). Monitor a cellular phone call and
observe the type of telephone and speed of the user. The probability model for this
experiment has the following information: P[$F$] = 0.5, P[$HF$] = 0.2, P[$MW$] = 0.1.
What is the sample space of the experiment? Calculate the following probabilities:
 
\textcolor{blue}{$S = \{HF, HW, MF, MW\}$}
\begin{enumerate}
    \item[a] \(P[W]\) = \textcolor{blue}{$P(F^\mathsf{c})= 1 - P(F)= 0.5$}
    \item[b] \(P[MF]\), \textcolor{blue}{$P(F) = P(HF)+P(MF),\quad P(MF)=P(F)-P(HF)=0.5-0.2=0.3$}
    \item[c] \(P[H]\)= \textcolor{blue}{$P(M^\mathsf{c})= 1- (P(MF)+P(MW)) = 1-0.3-0.1=0.6$}
\end{enumerate}

\section*{(3)}
Cellular phones perform handoffs as they move from cell to cell. During a
phone call, telephone either performs zero handoffs ($H_0$), one handoff ($H_1$), or more
than one handoff ($H_2$). In addition, each call is either long ($L$), if it lasts more than 3
minutes, or brief ($B$). The following table describes the possible types of calls.
\[
\begin{array}{c|ccc}
 & H_0 & H_1 & H_2 \\ \hline
L & 0.1 & 0.1 & 0.2 \\
B & 0.4 & 0.1 & 0.1 \\
\end{array}
\]
What is the probability $P[H_0]$ that a phone makes no handoffs? What is the probability
a call is brief? What is the probability that the call is long or there are at least two
handoffs?

\textcolor{blue}{$P[H_0]=P[H_0L]+P[H_0B]=0.4+0.1=0.5$}

\textcolor{blue}{$P[B]=P[H_0B]+P[H_1B]+P[H_2B]=0.4+0.1+0.1=0.6$}

\textcolor{blue}{$P[L \cup H_2]=P[L]+P[H_2]-P[L\cap H_2] = 0.4 + 0.3 -0.2 = 0.5$}

\newpage

\section*{(4)}
Monitor three consecutive phone calls going through a telephone switching
office. Classify each one as a voice call ($v$), if someone is speaking or a data call ($d$)
if the call is carrying a modem or fax signal. Your observation is a sequence of three
letters (\textit{each letter is either $v$ or $d$}). For example two voice calls followed by one data
call corresponds to $vvd$. Write the elements of the following sets:
Write the elements of:

\begin{enumerate}
    \item[a.] $A_1 = \{\text{first call is voice}\}=$ \textcolor{blue}{$\{vvv, vvd, vdv, vdd\}$}  
    \item[b.] $B_1 = \{\text{first call is data}\}=$  \textcolor{blue}{$\{dvv, dvd, ddv, ddd\}$}
    \item[c.] $A_2 = \{\text{second call is voice}\}=$  \textcolor{blue}{$\{vvv, vvd, dvv, dvd\}$}
    \item[d.] $B_2 = \{\text{second call is data}\}=$ \textcolor{blue}{$\{vdv, vdd, ddv, ddd\}$}
    \item[e.] $A_3 = \{\text{all calls the same}\}=$  \textcolor{blue}{$\{vvv, ddd\}$}
    \item[f.] $B_3 = \{\text{voice and data alternate}\}=$ \textcolor{blue}{$\{vdv, dvd\}$}
    \item[g.] $A_4 = \{\text{one or more voice calls}\}=$ \textcolor{blue}{$\{vvv, vvd, vdv, vdd, dvv, dvd, ddv\}$}
    \item[h.] $B_4 = \{\text{two or more data calls}\}=$ \textcolor{blue}{$\{vdd, dvd, ddv, ddd\}$} 
\end{enumerate}
For each pair of events $A_1$ and $B_1$, $A_2$ and $B_2$ and so on, please identify whether the
pair of events is either mutually exclusive or collectively exhaustive

\textcolor{blue}{$A_1$ and $B_1$ are collectively exhaustive and mutually exclusive.} 

\textcolor{blue}{$A_2$ and $B_2$ are collectively exhaustive and mutually exclusive.}

\textcolor{blue}{$A_3$ and $B_3$ are mutually exclusive.}

\textcolor{blue}{$A_4$ and $B_4$ are collectively exhaustive.}
\newpage 

\section*{(5)}
A student’s test score T is an integer between 0 and 100 corresponding to the
experimental outcomes $s0, \dots , s100$. $A$ score of 90 to 100 is an $A$, 80 to 89 is a $B$, 70
to 79 is a $C$, 60 to 69 is a $D$, and below 60 is a failing grade of $F$. Given all scores
between 51 and 100 are equally likely and a score of 50 or less cannot occur, find the
following probabilities.

\textcolor{blue}{Since all scores are equally likely, $P(Event) = \frac{\# of favorable outcomes}{Total \# of outcomes}$} 

\begin{enumerate}
    \item[a] \(P[\{s_{79}\}]\) = \textcolor{blue}{$\frac{1}{50} = 0.02$}
    \item[b] \(P[\{s_{100}\}]\) = \textcolor{blue}{$\frac{1}{50} = 0.02$}
    \item[c] \(P[A]\) = \textcolor{blue}{$\frac{11}{50} = 0.22$}
    \item[d] \(P[F]\) = \textcolor{blue}{$\frac{9}{50} = 0.18$}
    \item[e] \(P[T \ge 80]\) = \textcolor{blue}{$\frac{21}{50} = 0.42$}
    \item[f] \(P[T < 90]\) = \textcolor{blue}{$\frac{39}{50} = 0.78$}
    \item[g] \(P[\text{C grade or better}]\) = \textcolor{blue}{$\frac{31}{50} = 0.62$}
    \item[h] \(P[\text{student passes}]\) = \textcolor{blue}{$\frac{41}{50} = 0.82$}
\end{enumerate}

\section*{(6)}
Monitor a phone call. Classify the call as a voice call ($V$), if someone is
speaking; or a data call ($D$) if the call is carrying a modem or fax signal. Classify
the call as long ($L$) if the call lasts for more than three minutes; otherwise classify
the call as brief ($B$). Based on data collected by the telephone company, we use the
following probability model: $P[V ] = 0.7, P[L] = 0.6, P[V L] = 0.35.$ Find the following
probabilities:

\begin{enumerate}
    \item[a] \(P[DL]\) \textcolor{blue}{$ = P(L) -P(VL) = 0.6 -0.35 = 0.25$}
    \item[b] \(P[D \cup L]\) \textcolor{blue}{$ = P(D) + P(L) - P(D \cap L) = P(V^\mathsf{c})+ 0.6 - 0.25 = (1 - 0.7) + 0.6 - 0.25 = 0.65$}  
    \item[c] \(P[VB]\) \textcolor{blue}{$ = P(V)-P(VL) = 0.7 - 0.35 = 0.35$} 
    \item[d] \(P[V \cup L]\) \textcolor{blue}{$ = P(V) + P(L) - P(V \cap L) = 0.7 + 0.6 - 0.35 = 0.95$}
    \item[e] \(P[V \cup D]\) \textcolor{blue}{$ = P(S) = 1$ since V and D are collectively exhaustive.}
    \item[f] \(P[LB]\) \textcolor{blue}{$ = 0$ since L and B are mutally exclusive.}
\end{enumerate}

\end{document}
